% Frozen source-aware proof candidate.  Do not insert into the live chapter
% before production Rethlas verification.  No novelty or priority claim is
% intended.

\paragraph{Imported inputs.}
The proof uses the following primary-source statements with their stated
hypotheses.
\begin{enumerate}
\item Chambert--Loir--Ducros, Theorem~17.5.7(b), in the version of
25~July~2025, says that the strict integral residual functor on an
almost-good strict germ is essentially surjective already in the subcategory
whose underlying analytic bundle is trivial.  In rank one, every line bundle
on a model of the reduction of a good strict germ is therefore represented
by the residue of an integral PL germ.
\item Gubler--Jell--Rabinoff, \S\S6.2 and~6.6 and Lemma~6.5, identify the
Picard group of the reduction of a germ with the filtered colimit over its
models and provide common dominating models for finite collections of
residual line bundles.  Lemma~6.11 says that every model representing a nef
residual class is nef.  Lemma~6.18 gives linearity of residues, and
Proposition~7.1 identifies semipositivity at a point with nefness of the
residual class.  Remark~7.3(4) says that the semipositivity locus is open.
\item For a product principally polarized abelian variety, the Rosati map
identifies its real N\'eron--Severi space with the symmetric subspace of its
real endomorphism algebra and identifies the ample cone with the cone of
positive-definite matrices.  This is Prendergast-Smith, Theorem~4.3.  The
nef cone is its closure, the cone of positive-semidefinite matrices.
For \(S=E\times E\) with \(\operatorname{End}(E)=\mathbf Z\), the three
fibre and diagonal classes used below form an integral basis of
\(\operatorname{NS}(S)\); compare Bauer--Borntr\"ager, Lemma~3.1 and
Remark~3.4.
\item Chow's lemma \cite[Tag~02O2]{stacks-project} and resolution of
surface singularities \cite[Tag~0BGP]{stacks-project} give a smooth
projective surface dominating any proper surface model of the same function
field.  A proper birational morphism between smooth surfaces factors into
blowups in closed points \cite[Tag~0C5Q]{stacks-project}.  For such a
blowup \(b:Y'\to Y\), with exceptional curve \(E\), the sequence
\[
 0\longrightarrow\operatorname{Pic}(Y)
  \xrightarrow{b^*}\operatorname{Pic}(Y')
  \xrightarrow{\deg_E}\mathbf Z\longrightarrow0
\]
is exact \cite[Tag~0C5L]{stacks-project}.
\end{enumerate}

\begin{proof}[Proof of \cref{prop:p09-irrational-nef-ray-obstruction}]
Suppose, for a contradiction, that
\eqref{eq:p09-forbidden-positive-decomposition} holds.  Delete every term
with \(a_j=0\); thus all remaining \(a_j\) are strictly positive.  Taking
residues gives
\begin{equation}
\label{eq:p09-residual-positive-decomposition}
                 \alpha=\sum_{j=1}^{m}a_jL_{g_j}(x)
 \quad\text{in}\quad
 \operatorname{Pic}(P_{k(S)/k})\otimes_{\mathbf Z}\mathbf R.
\end{equation}
By the filtered common-model formalism
\cite[\S\S6.2, 6.6 and Lemma~6.5]{GJR26}, there is one model \(V\) of
the full Riemann--Zariski space \(P_{k(S)/k}\) dominating \(S\) and all
models involved in \eqref{eq:p09-residual-positive-decomposition}, far
enough that the equality itself is represented on \(V\).  Because it is a
model of the full Riemann--Zariski space, the morphism \(V\to S\) is proper.
Both models have function field \(k(S)\), so this morphism is birational.
Moreover, every integral line bundle on \(V\) representing
\(L_{g_j}(x)\) is nef by \cite[Lemma~6.11]{GJR26}.

Chow's lemma \cite[Tag~02O2]{stacks-project} and surface resolution
\cite[Tag~0BGP]{stacks-project} give a smooth projective surface \(Y\)
dominating \(V\).  Denote the composite morphism by
\[
                              f:Y\longrightarrow S.
\]
Pulling the represented equality to \(Y\), and then passing to numerical
classes, yields
\begin{equation}
\label{eq:p09-upstairs-nef-decomposition}
                       f^*[\alpha]=\sum_{j=1}^{m}a_j[M_j]
              \quad\text{in }N^1(Y)_{\mathbf R},
\end{equation}
where every \(M_j\) is an integral nef line bundle on \(Y\).

Let \(C\subset Y\) be an irreducible \(f\)-exceptional curve.  The
projection formula and \eqref{eq:p09-upstairs-nef-decomposition} give
\[
 0=f^*[\alpha]\cdot C
   =\sum_{j=1}^{m}a_j(M_j\cdot C).
\]
Every intersection number on the right is nonnegative, and every \(a_j\)
is positive.  Hence
\begin{equation}
\label{eq:p09-exceptional-orthogonality}
                              M_j\cdot C=0
\end{equation}
for every \(j\) and every \(f\)-exceptional curve \(C\).

By \cite[Tag~0C5Q]{stacks-project}, the morphism \(f\) factors as
\[
 Y=Y_r\xrightarrow{b_r}Y_{r-1}\xrightarrow{b_{r-1}}\cdots
   \xrightarrow{b_1}Y_0=S,
\]
where every \(b_\ell\) is the blowup of a closed point.  Fix \(j\), set
\(M_{j,r}=M_j\), and consider the exceptional curve \(E_r\subset Y_r\)
of \(b_r\).  It is \(f\)-exceptional, so
\eqref{eq:p09-exceptional-orthogonality} gives
\(\deg_{E_r}M_{j,r}=0\).  The exact Picard sequence for a point blowup
\cite[Tag~0C5L]{stacks-project} therefore supplies a line bundle
\(M_{j,r-1}\) on \(Y_{r-1}\) such that
\[
                              M_{j,r}\simeq b_r^*M_{j,r-1}.
\]
The line bundle \(M_{j,r-1}\) is nef: if \(D\subset Y_{r-1}\) is an
irreducible curve and \(\widetilde D\subset Y_r\) its strict transform,
then
\[
                     M_{j,r-1}\cdot D
                   = M_{j,r}\cdot\widetilde D\geq0.
\]
It also has degree zero on every curve contracted by
\(Y_{r-1}\to S\), since the strict transform of such a curve is
\(f\)-exceptional.  Iterating the same argument gives a nef line bundle
\(B_j=M_{j,0}\) on \(S\) with
\[
                         M_j\simeq f^*B_j.
\]
Consequently, for the integral nef class
\(\beta_j=c_1(B_j)\in N^1(S)_{\mathbf Z}\), one has
\[
                         [M_j]=f^*\beta_j
                  \quad\text{in }N^1(Y)_{\mathbf R}.
\]

Substituting into \eqref{eq:p09-upstairs-nef-decomposition} and using the
injectivity of \(f^*:N^1(S)_{\mathbf R}\to N^1(Y)_{\mathbf R}\), we obtain
\begin{equation}
\label{eq:p09-downstairs-nef-decomposition}
                              [\alpha]=\sum_{j=1}^{m}a_j\beta_j
                       \quad\text{in }N^1(S)_{\mathbf R}.
\end{equation}
Each summand is nef.  Since \([\alpha]\) spans an extremal nef ray,
every nonzero \(\beta_j\) belongs to \(\mathbf R_{\geq0}[\alpha]\).  At
least one \(\beta_j\) is nonzero because \([\alpha]\ne0\).  This contradicts
\eqref{eq:p09-irrational-extremal-ray}, since
\(\beta_j\in N^1(S)_{\mathbf Z}\subset N^1(S)_{\mathbf Q}\).
Therefore the positive decomposition cannot exist.

For the last assertion, write the residual class as a finite sum
\[
                         \alpha=\sum_{i=1}^{q}\lambda_i[L_i],
 \qquad
             \lambda_i\in\Lambda,\quad L_i\in\operatorname{Pic}(S).
\]
Replace \(X\) by a strictly affinoid neighbourhood of \(x\).  It is
separated, and its good germ at \(x\) is almost good.
The strict residual lifting theorem
\cite[Theorem~17.5.7(b)]{CLD25} supplies a strictly integral PL metric on
an analytically trivial line bundle for each residual bundle \([L_i]\).
Choose a local trivializing frame and let \(u_i\) be the negative logarithm
of its norm.  Compatibility of the strict residual functor with the residue
line-bundle convention, exactly as in the realization step in the proof of
\cite[Proposition~7.10]{GJR26}, gives
\[
                              L_{u_i}(x)=[L_i].
\]
After shrinking once so that all these germs are represented by functions,
put
\[
                              h_\alpha=\sum_{i=1}^{q}\lambda_i u_i.
\]
This is a \(\Lambda\)-PL germ.  Residue linearity
\cite[Lemma~6.18]{GJR26} gives \(L_{h_\alpha}(x)=\alpha\), and the nefness
of \(\alpha\), together with the residue criterion
\cite[Proposition~7.1]{GJR26}, shows that \(h_\alpha\) is semipositive at
\(x\).
\end{proof}

\begin{proof}[Proof of
\cref{cor:p09-good-reduction-positive-cone-counterexample}]
Set \(\tau=\pi\mathrm i\).  We first check that \(E\) has no complex
multiplication.  If multiplication by \(z\in\mathbf C\) preserves
\(\mathbf Z+\tau\mathbf Z\), then for some integers \(a,b,c,d\) one has
\[
                              z=a+b\tau,
 \qquad
                              z\tau=c+d\tau.
\]
Thus
\[
                       b\tau^2+(a-d)\tau-c=0.
\]
The number \(\tau\) is transcendental, so \(b=c=0\) and \(a=d\).  Hence
\(z=a\in\mathbf Z\), proving
\[
                              \operatorname{End}(E)=\mathbf Z.
\]

Let \(p_1,p_2:S\to E\) be the projections, and put
\[
 H_1=p_1^*(0)=\{0\}\times E,
 \qquad
 H_2=p_2^*(0)=E\times\{0\},
 \qquad
 \Delta=(p_1-p_2)^*(0).
\]
The classes \(H_1,H_2,\Delta\) form a \(\mathbf Z\)-basis of
\(\operatorname{NS}(S)\) by
\cite[Remark~3.4]{BauerBorntraeger18}.  For completeness, the Rosati
identification is
\[
\rho:\operatorname{NS}(S)\simeq
 \left\{
   \begin{pmatrix}r&s\\s&u\end{pmatrix}:r,s,u\in\mathbf Z
 \right\}
\]
for the product principal polarization \(\Theta=H_1+H_2\), where
\(\rho([D])=\phi_\Theta^{-1}\circ\phi_D\).  Under this identification,
\[
 H_1\longmapsto
 \begin{pmatrix}1&0\\0&0\end{pmatrix},
 \qquad
 H_2\longmapsto
 \begin{pmatrix}0&0\\0&1\end{pmatrix},
 \qquad
 \Delta\longmapsto
 \begin{pmatrix}1&-1\\-1&1\end{pmatrix}.
\]
Indeed, if \(\Theta_E=(0)\) is the principal polarization on \(E\), the
functorial formula for the polarization homomorphism of \(q^*\Theta_E\)
gives the symmetric matrix \(q^{\mathsf T}q\).  Taking
\(q=p_1,p_2,p_1-p_2\) gives exactly the three displayed matrices and fixes
the sign of the off-diagonal entries.
The ample cone corresponds to the positive-definite cone by
\cite[Theorem~4.3]{PrendergastSmith12}, and therefore the nef cone
corresponds to the positive-semidefinite cone.

Put \(s=\sqrt2\) and consider the Picard class
\[
 \widehat\xi=(2+s)[\mathcal O_S(H_1)]
        +(1+s)[\mathcal O_S(H_2)]-s[\mathcal O_S(\Delta)]
        \in\operatorname{Pic}(S)\otimes_{\mathbf Z}\Lambda
\]
and its numerical image
\begin{equation}
\label{eq:p09-explicit-irrational-nef-class}
               \xi=(2+s)H_1+(1+s)H_2-s\Delta
                 \in\operatorname{NS}(S)\otimes_{\mathbf Z}\Lambda.
\end{equation}
Its Rosati matrix is
\begin{equation}
\label{eq:p09-rank-one-rosati-matrix}
 Q=
 \begin{pmatrix}2&s\\s&1\end{pmatrix}
 =
 \begin{pmatrix}s\\1\end{pmatrix}
 \begin{pmatrix}s&1\end{pmatrix}.
\end{equation}
Thus \(\xi\) is nonzero and nef.  Its ray is extremal.  Indeed, if
\(Q=A+B\) with \(A,B\) positive semidefinite and if \(w\in\ker Q\), then
\[
                    0=w^{\mathsf T}Aw+w^{\mathsf T}Bw.
\]
Both summands are nonnegative, so both vanish.  For a positive-semidefinite
matrix, \(w^{\mathsf T}Aw=0\) implies \(Aw=0\); hence
\(\ker Q\subset\ker A\cap\ker B\).  Since \(Q\) has rank one, the images
of \(A\) and \(B\) lie in the one-dimensional image of \(Q\), and
\(A,B\) are nonnegative scalar multiples of \(Q\).

The ray is irrational.  If \(r\xi\in N^1(S)_{\mathbf Q}\), the difference
between its \(H_1\)- and \(H_2\)-coefficients is \(r\), whereas its
\(\Delta\)-coefficient is \(-r\sqrt2\).  Both are rational only when
\(r=0\).  Thus
\[
                   \mathbf R_{\geq0}\xi\cap N^1(S)_{\mathbf Q}=\{0\}.
\]

For the divisorial point \(x\) of the constant good-reduction model, the
residue of its completed residue field is the function field of the closure
of its reduction, namely
\[
 \widetilde{\mathscr H(x)}=\mathbf C(S).
\]
The construction in \cite[\S6.13]{GJR26} identifies the reduction of the
germ with
\[
 \operatorname{red}(X,x)=P_{\mathbf C(S)/\mathbf C}\{S\}.
\]
Since \(S\) is proper, the valuative criterion gives
\(P_{\mathbf C(S)/\mathbf C}\{S\}=P_{\mathbf C(S)/\mathbf C}\).
The last assertion of
\cref{prop:p09-irrational-nef-ray-obstruction}, applied to
\(\alpha=\widehat\xi\),
gives a semipositive \(\Lambda\)-PL germ with residue \(\widehat\xi\); the first
assertion excludes every nonnegative-real decomposition of this germ into
semipositive integral PL germs.  By the openness of its
semipositivity locus \cite[Remark~7.3(4)]{GJR26}, the germ is represented by
a semipositive function on one strict analytic neighbourhood \(U\) of
\(x\), as asserted.
\end{proof}

\begin{remark}
\label{rem:p09-positive-coefficients-essential}
The sign restriction in
\eqref{eq:p09-forbidden-positive-decomposition} is essential.  Each of
\(H_1,H_2,\Delta\) is nef, so their integral PL lifts may be chosen
semipositive at \(x\).  Formula
\eqref{eq:p09-explicit-irrational-nef-class} therefore writes the same
counterexample as a signed \(\Lambda\)-linear combination of semipositive
integral PL germs.  With negative coefficients, intersections against
exceptional curves can cancel and the descent argument above no longer
applies.
\end{remark}
